% ===================================================== BÖLÜM 6
\section{Bilgisayar Kurulumu ve Kodlar (Hepsi Test Edildi)}
\label{sec:kod}

\subsection{En kolay yol: Google Colab}
Bilgisayarına hiçbir şey kurmadan çalışmak için:
\begin{enumerate}[leftmargin=*]
\item \texttt{colab.research.google.com} $\rightarrow$ “Yeni not defteri”.
\item İlk hücreye yaz ve çalıştır (Shift+Enter):
\begin{lstlisting}[style=kabuk]
!pip install -q astropy astroquery batman-package emcee
\end{lstlisting}
\item Sol menüdeki klasör simgesinden aşağıdaki \texttt{.py} dosyalarını yükle (ya da bir hücreye
\texttt{\%\%writefile oc\_diyagrami.py} yazıp altına kodu yapıştır).
\item Komutları başına \texttt{!} koyarak çalıştır: \texttt{!python oc\_diyagrami.py "WASP-12 b" WASP-012}.
\end{enumerate}
\textbf{Kendi bilgisayarında} (Windows/Mac/Linux): Anaconda kur, “Anaconda Prompt”ta şunu yaz:
\begin{lstlisting}[style=kabuk]
pip install numpy scipy matplotlib astropy astroquery batman-package emcee
\end{lstlisting}
Windows'ta
\texttt{batman-package} derlenemezse Colab kullan.

\subsection{Klasör düzeni}
\begin{lstlisting}[style=kabuk]
tubitak_2209a/
  kod/        oc_diyagrami.py  tess_zaman.py  rv_analiz.py  fizik_testleri.py  zaman_donusum.py  secim_tablosu.py
  veri/       tess_WASP-12.csv ...  (kendi olctugun TESS zamanlari; zamanlar_*.csv: O-C tablolari)
  sekiller/   OC_*.png  RV_*.png
\end{lstlisting}
\textbf{Kural}: İndirilen ham veriye elle dokunma; her dönüşüm koddan yapılsın. Böylece her sonuç yeniden üretilebilir.

\subsection{Kod 1: \texttt{oc\_diyagrami.py} — O$-$C diyagramı, $\dot P$, BIC, LOOCV}
Ne yapar? (1) Ivshina \& Winn ve Wang vd. tablolarını VizieR'den indirir; (2) zaman sistemini kontrol eder;
(3) aynı geçişi iki kez saymaz; (4) istersen kendi TESS zamanlarını ekler; (5) epokları hesaplar; (6) 5$\sigma$ kırpma
yapar; (7) doğrusal ve kuadratik modeli uydurur; (8) $\dot P$ ve $\Delta$BIC'i yazar (ayrıca bilgi amaçlı LOOCV sonucunu da yazar; projede kullanılmaz); (9) O$-$C
diyagramını \texttt{OC\_<gezegen>.png} olarak, tabloyu \texttt{zamanlar\_<gezegen>.csv} olarak kaydeder.

%%KOD:oc_diyagrami.py%%

\textbf{Çalıştırma ve beklenen çıktı} (Kasım 2026, yalnız katalog verisi):
\begin{lstlisting}[style=kabuk]
python oc_diyagrami.py "WASP-12 b" WASP-012
WASP-12 b: n=189 (TESS 82), dP/dt = -28.64 +- 0.98 ms/yil, DBIC = 854.6, LOOCV: 0 noktada kosul bozuluyor -> GECTI
python oc_diyagrami.py "WASP-12 b" WASP-012 tess_WASP-12.csv      # kendi TESS zamanlarinla
WASP-12 b: n=227 (TESS 120), dP/dt = -31.05 +- 0.82 ms/yil, DBIC = 1422.4, LOOCV: ... -> GECTI
\end{lstlisting}
Ivshina \& Winn tablosundaki sistem adları: \texttt{WASP-012}, \texttt{WASP-004}, \texttt{HAT-P-07}, \texttt{TrES-1}
(diğerleri: \texttt{CoRoT-02}, \texttt{KELT-16}, \texttt{TrES-5}, \texttt{WASP-010}, \texttt{WASP-022},
\texttt{WASP-045}, \texttt{XO-4}, \texttt{HAT-P-37}, \texttt{WASP-001}, \texttt{WASP-046}).

\textbf{Kod nasıl okunur? (satır satır mantık)}
\begin{itemize}[leftmargin=*]
\item \texttt{vizier()}: Bir web adresi kurar, VizieR'den sekmeyle ayrılmış metin indirir, \texttt{\#} ile başlayan
      açıklama satırlarını atar, ilk satırı sütun adı yapar, birim ve çizgi satırlarını atlar.
\item \texttt{zamanlari\_topla()}: “O”ları toplar. \texttt{n\_TTmid} sütunu \texttt{BJD\_TDB} ya da \texttt{BJD\_TT}
      değilse uyarı verir. \texttt{BELIRSIZ\_BJD\_DAHIL = False} yaparak belirsiz “BJD” zamanlarını dışarıda bırakıp
      sonucun değişip değişmediğine bakabilirsin (HAT-P-7 dersi).
\item \texttt{uydur()}: Ağırlıklı en küçük kareler; \texttt{derece=1} doğrusal, \texttt{2} kuadratik, \texttt{3} kübik.
\item \texttt{analiz()}: $\dot P = 2c_2/P$ dönüşümü ve $\Delta$BIC.
\item \texttt{sigma\_kirp()}: 5$\sigma$ dışındaki noktaları yinelemeli atar.
\item \texttt{loocv()}: Bölüm~\ref{sec:loocv}'teki algoritmanın birebir karşılığı.
\item Ana blok: “C” için Wang Tablo 2'deki $T_0$, $P$ ile epokları bulur; çizimde O$-$C'yi kendi doğrusal
      uydurmasına göre hesaplar.
\end{itemize}

\subsection{Kod 2: \texttt{tess\_zaman.py} — TESS ışık eğrisinden orta geçiş zamanı (O) ölçmek}
Girdi: hedef adı, RA, Dec (derece), $P$, $T_0$ ve sektör listesi. Hangi sektörlerin var olduğunu MAST'ta
(\texttt{mast.stsci.edu}, “TESS” filtresi) ya da aşağıdaki komutla görebilirsin. Wang vd.'nin kullanmadığı
(sonraki) sektörler, Kasım 2026 itibarıyla:

\begin{center}
\small
\begin{tabularx}{\textwidth}{l >{\raggedright\arraybackslash}X l >{\raggedright\arraybackslash}X}
\toprule
Hedef & Tüm SPOC sektörleri & Wang'ın son TESS verisi & Yeni sektörler (bu projede ölçülen)\\
\midrule
WASP-12 & 20, 43–46, 71, 72 & Aralık 2021 & 71, 72 (38 geçiş)\\
WASP-4  & 2, 28, 29, 36, 69, 96, 103, 105, 106 & Eylül 2020 & 36, 69, 96, 103, 105, 106 (75 geçiş)\\
HAT-P-7 & 14–16, 19, 40, 41, 54, 55, 74, 75, 78, 81, 82, 86 & Ağustos 2022 & 74, 75, 78, 81, 82, 86 (40 geçiş)\\
TrES-1  & 14, 40, 41, 53–55, 74, 78, 80, 81, 86 & Ağustos 2022 & 74, 78, 80, 81, 86 (21 geçiş)\\
\bottomrule
\end{tabularx}
\end{center}

%%KOD:tess_zaman.py%%

\begin{lstlisting}[style=kabuk]
python tess_zaman.py WASP-12 97.6366450 29.6722662 1.091418901 2457607.519305 71 72
Tur 1: Rp/R* = 0.1135, a/R* = 3.07, i = 84.52
Tur 2: Rp/R* = 0.1131, a/R* = 3.12, i = 85.72
Tc = 2460235.652614 BJD_TDB  +-  43.9 s
...   (38 gecis; sonuc tess_WASP-12.csv dosyasina yazilir)
python tess_zaman.py HAT-P-7 292.2471900 47.9695429 2.204737 2454954.358572 74 75 78 81 82 86
python tess_zaman.py TrES-1 286.0408755 36.6325360 3.03007 2453186.806341 74 78 80 81 86
python tess_zaman.py WASP-4 353.5629149 -42.0617794 1.338231 2455804.515679 36 69 96 103 105 106
\end{lstlisting}
Bir sektörün indirilmesi ve ölçülmesi birkaç dakika sürer (MCMC). \textbf{Sağlama}: “Tur 2” satırındaki $R_p/R_\star$ ve
$a/R_\star$ literatüre yakın olmalı (WASP-12: 0,117 / 3,0; HAT-P-7: 0,078 / 4,2; TrES-1: 0,137 / 10,5; WASP-4: 0,152 / 5,5).
Çok farklıysa (ör. derinlik yarıya inmişse) efemeris kaymıştır ya da başka hedefin dosyası karışmıştır.

\subsection{Kod 3: \texttt{rv\_analiz.py} — Arşiv RV'sinden Rømer testi}
Bölüm~\ref{sec:rv}'te ayrıntılı. Arşivdeki hazır RV'leri (Knutson vd. 2014; Bouma vd. 2020 ya da senin derlediğin bir
CSV) indirir; her alete ayrı sıfır noktası, gezegenin dairesel yörünge sinyali ve doğrusal bir eğilim uydurur;
eğilimden Rømer'in öngördüğü $\dot P$'yi hesaplar ve O$-$C'den bulduğunla karşılaştırır.

%%KOD:rv_analiz.py%%

\subsection{Kod 4: \texttt{fizik\_testleri.py} — RV gerektirmeyen fizik testleri}
Gelgit için gereken Q$'_\star$, Roche oranı (kütle kaybı), apsis genliği üst sınırı, eğriliğin sabitliği (kübik model)
ve periyodiklik (Lomb–Scargle periyodogramı).

%%KOD:fizik_testleri.py%%

\subsection{Kod 5: \texttt{zaman\_donusum.py} — HJD/BJD$_\mathrm{UTC}$ $\rightarrow$ BJD$_\mathrm{TDB}$}
ETD/VarAstro'dan ya da eski makalelerden gelen zamanlar için.

%%KOD:zaman_donusum.py%%

\begin{lstlisting}[style=kabuk]
python zaman_donusum.py 2456000.5 HJD_UTC 97.6366450 29.6722662
HJD_UTC 2456000.500000 -> BJD_TDB 2456000.500759   (fark +65.6 s)
\end{lstlisting}

\subsection{Kod 6: \texttt{dace\_rv\_indir.py} — DACE arşivinden RV indirmek}
Bir yıldızın DACE'deki yayımlanmış RV'lerini indirir, geçiş sırasında alınanları (Rossiter–McLaughlin etkisi, yörünge
sinyalini bozar) ve 3'ten az ölçümü olan aletleri atar, \texttt{rv\_analiz.py}'nin okuduğu CSV'yi yazar.

%%KOD:dace_rv_indir.py%%

\subsection{Sık karşılaşılan hatalar}
\begin{itemize}[leftmargin=*]
\item \textbf{VizieR açılmıyor / bağlantı reddediliyor}: Kodda \texttt{VIZIER} adresini
      \url{https://vizier.cds.unistra.fr/viz-bin/asu-tsv} ile değiştir (ya da tersi).
\item \textbf{\texttt{IndexError} (\texttt{[0]})}: Gezegen adı yanlış yazılmış. Wang tablosunda \texttt{"WASP-12 b"}
      (boşluklu), I\&W'de \texttt{WASP-012}.
\item \textbf{O$-$C'de bir nokta $\pm P/2$ kadar uzakta}: Epok bir tur kaymış ya da ikincil tutulma zamanı karışmış.
\item \textbf{Bir kaynağın bütün noktaları $\sim$1 dakika kaymış}: Zaman sistemi (UTC/TDB ya da HJD) hatası.
      \texttt{zaman\_donusum.py} ile düzelt.
\item \textbf{$\Delta$BIC çok büyük ama O$-$C'de parabol görünmüyor}: Hata çubukları gerçekçi değil (çok küçük).
      İndirgenmiş $\chi^2$ ($\chi^2/(n-3)$) 1'den çok büyükse, hataları $\sqrt{\chi^2_\nu}$ ile büyütüp yeniden dene.
\end{itemize}

% ===================================================== BÖLÜM 7
\section{Dönem Değişiminin 6 Olası Sebebi}
\label{sec:etkiler}

Hocanın listesindeki altı etkinin her biri için aynı sırayla gideceğiz: \textbf{Ne?} (basitçe),
\textbf{O$-$C'deki izi}, \textbf{formül}, \textbf{bizim hedeflerde büyüklüğü}, \textbf{nasıl test edilir}.
Başlamadan önce yedinci ve en sık sebebi unutma: \textbf{veri hatası} (yanlış zaman sistemi, gerçekçi olmayan
küçük hata, yanlış epok). Diğer altısına bakmadan önce Bölüm~\ref{sec:oc}'teki temizlik (zaman sistemi, epok, aykırı nokta) yapılmış olmalı.

\subsection{Etki 1: Gelgit etkileşimi (tidal interaction / orbital decay)}
\textbf{Ne?} Gezegen yıldızda (Ay'ın Dünya okyanuslarında yaptığı gibi) iki gelgit kabarıklığı oluşturur. Sıcak
Jüpiterlerin çoğunda yıldız kendi ekseni etrafında, gezegenin yörüngede dolanmasından \emph{yavaş} döner. Bu yüzden
kabarıklık gezegenin gerisinde kalır ve onu geriye doğru çeker: gezegen enerji kaybeder, yavaş yavaş içeri doğru
sarmallanır, \textbf{yörünge dönemi kısalır}. Sonunda gezegen yıldız tarafından yutulur.

\textbf{O$-$C izi}: Aşağı bükülen, düzgün bir \textbf{parabol} (Şekil~\ref{fig:katalog}c). $\dot P$ on yıllar boyunca
sabittir (kübik terim gerekmez). İşaret: \textbf{negatif} (yıldız hızlı dönüyorsa pozitif olabilir; ama hedeflerimizin
hiçbiri öyle değil: Wang vd. 2024, Bölüm 5.3).

\textbf{Formül} (sabit Q$'_\star$ ile; ör. Patra vd. 2017, Yee vd. 2020):
\[
\dot P = -\frac{27\pi}{2\,Q'_\star}\left(\frac{M_p}{M_\star}\right)\left(\frac{R_\star}{a}\right)^5
\qquad\Longleftrightarrow\qquad
Q'_\star = -\frac{27\pi}{2}\left(\frac{M_p}{M_\star}\right)\left(\frac{R_\star}{a}\right)^5\frac{1}{\dot P}.
\]
$Q'_\star$ (değiştirilmiş gelgit kalite faktörü), yıldızın gelgit enerjisini ne kadar verimli soğurduğunun ölçüsüdür.
\textbf{Küçük Q$'_\star$ = çok soğurma = hızlı bozulma.} Kuram ve gözlemler $Q'_\star \sim 10^5$–$10^7$ bekler
(Penev vd. 2018). $(R_\star/a)^5$ çok güçlü bir terimdir: yıldıza yakın ve büyük yıldızlı sistemler en hızlı bozulur.

\textbf{Bizim hedeflerde} (\texttt{fizik\_testleri.py}):
\begin{center}
\small
\begin{tabular}{l r r l}
\toprule
Hedef & $\dot P$ (ms/yıl) & Gereken Q$'_\star$ & Yorum\\
\midrule
WASP-12\,b & $-31{,}05\pm0{,}82$ & $1{,}9\times10^5$ & Makul. Gelgit mümkün.\\
WASP-4\,b  & $-5{,}11\pm0{,}30$  & $7{,}7\times10^4$ & Sınırda (beklenenden biraz küçük); mümkün ama zorlayıcı.\\
HAT-P-7\,b & $+28{,}14\pm3{,}26$ & -- & Dönem \textbf{uzuyor}: gelgit (yavaş dönen yıldız) bunu yapamaz.\\
TrES-1\,b  & $-12{,}70\pm1{,}90$ & $8{,}4\times10^2$ & Kuramdan $\sim$100 kat küçük: gelgit fiziksel olarak çok zor.\\
\bottomrule
\end{tabular}
\end{center}

\textbf{Nasıl test edilir?} (1) Q$'_\star$ makul aralıkta mı? (2) RV'de eğilim yok mu (gelgit sistemin kütle merkezini
ivmelendirmez)? (3) $\dot P$ zaman içinde sabit mi (kübik terim gereksiz mi)? (4) Örtülme zamanları geçişlerle
\emph{aynı yönde} mi kayıyor? Dördü de “evet” ise gelgit en olası sebeptir.

\subsection{Etki 2: Apsis hareketi (apsidal motion / precession)}
\textbf{Ne?} Yörünge tam çember değil de çok hafif elips ise ($e>0$), elipsin en yakın noktası (enberi) zamanla
yörünge düzleminde döner. Bu dönme, geçişin elips üzerinde gerçekleştiği yeri değiştirir; geçişler bir erken bir geç
gelir.

\textbf{O$-$C izi}: Dönemi yıllar–on yıllar olan bir \textbf{sinüs}. Gözlem süresi bu dönemden kısaysa sinüsün bir
parçası \textbf{parabol gibi} görünür; gelgitle karışabilir. \textbf{Ayırt edici imza}: \textbf{örtülme} (gezegenin
yıldızın arkasına geçtiği ikincil tutulma) zamanlarının O$-$C'si geçişlerinkiyle \textbf{ters yönde} salınır.
Gelgit ve Rømer'de ikisi \emph{aynı yönde} kayar.

\textbf{Formül} (Patra vd. 2017):
\[
T_\text{geçiş}(E) = T_0 + P_s E - \frac{e P_a}{\pi}\cos\omega(E),\quad
T_\text{örtülme}(E) = T_0 + \frac{P_a}{2} + P_s E + \frac{e P_a}{\pi}\cos\omega(E),\quad
\omega(E) = \omega_0 + \frac{d\omega}{dE}E .
\]
\textbf{Genlik} $\approx eP/\pi$: $e = 0{,}01$ ve $P = 1$ gün için $\sim$275 s. Çok küçük dışmerkezlik bile dakika
mertebesinde salınım yaratabilir.

\textbf{Bizim hedeflerde}: Dört hedefte de NASA arşivindeki $e$ ya sıfır ya da sıfırla uyumlu. Yine de üst sınırlarla
genlik, gözlenen sinyalden büyük olabilir (WASP-12: $\lesssim$1600 s'ye karşı gözlenen 415 s), yani \emph{yalnızca}
genliğe bakarak apsisi dışlayamayız. \textbf{En güçlü test WASP-12 için arşivde mevcut}: Yee vd. (2020) Tablo 5'te 19
örtülme zamanı var. Kendi hesabımda bu örtülmeler, geçişlerden çıkan kuadratik bozulma eğrisine $\chi^2/n = 1{,}07$
ile uyuyor, yani geçişlerle \textbf{aynı yönde} kayıyorlar. Kalan $\sim$34 s'lik sabit fark beklenen bir
şeydir: ışığın yörüngenin bir ucundan öbürüne gitme süresi $2a/c \approx 23$ s. \textbf{Sonuç: WASP-12'de apsis
hareketi gerekmiyor} (Yee vd. 2020 ile uyumlu). Diğer hedeflerde, kübik/sinüs testi yapılır ve TESS'te örtülme varsa
(HAT-P-7 parlak ve sıcak; örtülmesi Kepler'de ölçülmüştür) literatürden alınır.

\subsection{Etki 3: Rømer etkisi (ışık-yolu etkisi, LTTE)}
\textbf{Ne?} Sistemde uzakta (birkaç AB ötede) ikinci bir ağır cisim (dev gezegen, kahverengi cüce ya da sönük yıldız)
varsa, yıldız-gezegen çifti bu cisimle birlikte ortak kütle merkezi etrafında yavaşça döner. Çift bize yaklaşırken
geçişin ışığı bize daha kısa yoldan gelir (geçişler erken görünür); uzaklaşırken geç görünür. \textbf{Gezegenin yörüngesi
hiç değişmez}; değişen, ışığın bize gelme süresidir. Ole Rømer 1676'da Jüpiter'in uydularının tutulmalarında bu etkiyle
ışık hızını ölçmüştü.

\textbf{O$-$C izi}: Eşlikçinin yörünge dönemi gözlem süresiyle karşılaştırılabilirse \textbf{sinüs}; çok uzunsa
yalnızca sinüsün bir parçası görünür ve bu parça \textbf{parabol} gibidir. Bize doğru ivme $\Rightarrow$ dönem kısalıyor
gibi görünür (aşağı parabol); bizden uzağa ivme $\Rightarrow$ uzuyor gibi (yukarı parabol). Uzun zaman tabanında ivme
değişirse $\dot P$ de değişir (kübik terim gerekir).

\textbf{Formüller} (Bouma vd. 2020, Denk. 1; Winn \& Stefánsson 2025):
\[
\boxed{\;\dot P_\text{Rømer} = \frac{P\,\dot v_r}{c}\;}
\qquad\qquad
\text{LTTE genliği: } A = \frac{a_\star \sin i_c}{c},\quad a_\star = a_c\frac{M_c}{M_\star+M_c}.
\]
$\dot v_r$: yıldızın bakış doğrultusundaki ivmesi (RV eğilimi), $c$: ışık hızı. Örnek: WASP-4 için
$\dot v_r = -0{,}0422$ m/s/gün, $P = 1{,}338$ gün $\Rightarrow$ $\dot P_\text{Rømer} = -5{,}94$\msyil{} (kodum aynı sayıyı
veriyor). 8 AB'de 7 M$_\mathrm{J}$'lik bir eşlikçi için $A \approx 30$ s.

\textbf{Nasıl test edilir?} \textbf{Arşiv RV'si} (Bölüm~\ref{sec:rv}). Yalnız O$-$C ile gelgitten ayrılamaz.

\subsection{Etki 4: Applegate etkisi (manyetik etkinlik)}
\textbf{Ne?} Güneş gibi manyetik olarak etkin yıldızlar 11 yıllık çevrimler gibi etkinlik döngüleri yaşar. Applegate
(1992)'ye göre bu döngü yıldızın iç kütle dağılımını, yani basıklığını (kütle dört-kutup momentini, $Q$) değiştirir.
Gezegen biraz farklı bir çekim alanı hisseder ve yörünge dönemi \textbf{döngüyle birlikte bir uzar bir kısalır}.
Başlangıçta örten çift yıldızlar için önerildi; Watson \& Marsh (2010) sıcak Jüpiterlere uyguladı.

\textbf{O$-$C izi}: \textbf{Yarı periyodik} salınım. Dönemi yıldızın etkinlik döngüsü kadar (5–50 yıl). Kısa veri
parçasında parabol gibi görünebilir.

\textbf{Formül} (Applegate 1992):
\[
\frac{\Delta P}{P} = -9\left(\frac{R_\star}{a}\right)^2\frac{\Delta Q}{M_\star R_\star^2},
\qquad
\text{O$-$C genliği} = \frac{\Delta P}{P}\cdot\frac{P_\text{mod}}{2\pi}.
\]
Watson \& Marsh (2010): sıcak Jüpiterlerde genlik 11 yıllık döngü için birkaç saniye, 50 yıllık döngü için birkaç
dakikadır.

\textbf{Bizim hedeflerde}: Etkin yıldızlar geç G ve K türüdür. WASP-4 (G7V, 5400 K) ve TrES-1 (K0V, 5200 K) aday;
WASP-12 ve HAT-P-7 (F türü, 6300 K) zayıf etkinliklidir. Bouma vd. (2020) Tablo 2'de WASP-4 için HIRES S-indeksi
(kromosferik etkinlik ölçüsü) var.

\textbf{Nasıl test edilir? (basit sürüm)} (1) Yıldız etkin bir G/K cücesi mi? (2) Kuadratik model artıklarında
periyodiklik var mı (Lomb–Scargle, yanlış alarm olasılığı $<0{,}01$)? (3) Eğrilik sabit mi (kübik test)?
(4) Periyot 5–50 yıl aralığında mı? \textbf{Dikkat:} 1 yıla yakın periyotlar çoğunlukla gözlem mevsimlerinin izidir
(“alias”), gerçek sinyal sayılmaz. Veri birkaç döngüyü kapsamadıkça Applegate ancak “olası” ya da “dışlanamadı” olarak
raporlanır. Bu dürüst ve kabul edilen bir sonuçtur.

\subsection{Etki 5: Üçüncü cisim (dinamik etki)}
\textbf{Ne?} Hocanın listesinde Rømer ile “3. cisim” ayrı yazılmış. Fark şudur:
\begin{itemize}[leftmargin=*]
\item \textbf{Rømer}: Üçüncü cisim \emph{bütün sistemi} sürükler; gezegenin yörüngesi değişmez, yalnız ışık yolu değişir.
      \emph{Geometrik} bir etkidir.
\item \textbf{Dinamik 3. cisim}: Üçüncü cisim (özellikle yakındaki bir gezegen) \emph{gezegenin yörüngesini} kütle
      çekimiyle gerçekten bozar. Örnek: rezonanstaki komşu gezegen kısa dönemli geçiş zamanı değişimleri (TTV)
      üretir (WASP-47 sisteminde olduğu gibi). Ya da gizli bir yakın eşlikçi yörüngeyi presesyona zorlar
      (Hagey vd. 2025, TrES-1 için böyle bir eşlikçi önerdi).
\end{itemize}
\textbf{O$-$C izi}: Haftalar–aylar ölçeğinde kısa dönemli dalgalanmalar (rezonans TTV'si) ya da zorlanmış apsis
hareketi (sinüs).

\textbf{Nasıl test edilir?} (1) Kuadratik artıklarda Lomb–Scargle ile kısa dönemli sinyal ara (TESS'in sürekli 27 günlük
sektörleri bunun için idealdir). (2) Literatürde sistemde başka gezegen/eşlikçi bildirilmiş mi (RV, görüntüleme)?
(3) RV artıklarında başka sinyal var mı? Bulunmazsa: “dinamik 3. cisim etkisi için kanıt yok”.

\subsection{Etki 6: Kütle kaybı (mass loss)}
\textbf{Ne?} Yıldızına çok yakın gezegenler ısınıp şişer; atmosferleri buharlaşabilir ya da gezegen Roche lobunu
doldurup yıldıza madde aktarabilir. Kütle kaybı açısal momentum alışverişiyle yörüngeyi değiştirebilir (çoğunlukla
dışa doğru, yani dönem uzar; koşullara bağlı; Valsecchi vd. 2015, Jackson vd. 2016).

\textbf{O$-$C izi}: Tekdüze, çok yavaş bir parabol. Etkisi yalnızca gezegen Roche sınırına çok yakınsa önemlidir.

\textbf{Formül} (Roche sınırı): $a_\text{Roche} \approx 2{,}16\,R_p\,(M_\star/M_p)^{1/3}$.
$a/a_\text{Roche} \approx 1$ ise önemli olabilir; $\gtrsim 2$ ise ihmal edilir.

\textbf{Bizim hedeflerde}: WASP-12\,b \textbf{1,18} (Roche'a çok yakın; WASP-12\,b'nin madde kaybettiği biliniyor),
WASP-4\,b 1,76, HAT-P-7\,b 2,54, TrES-1\,b 3,09. Kütle kaybı WASP-12 dışında ihmal edilebilir. WASP-12'de de gözlenen
hızlı \emph{kısalmayı} açıklayamaz (kütle kaybı çoğunlukla dönemi uzatır). Wang vd. (2024) de kütle kaybının dönem
artışlarını açıklamadığını not eder.

\subsection{Özet tablo: hangi etki nasıl bir iz bırakır?}
\begin{center}
\small
\begin{tabularx}{\textwidth}{l >{\raggedright\arraybackslash}X >{\raggedright\arraybackslash}X >{\raggedright\arraybackslash}X >{\raggedright\arraybackslash}X}
\toprule
Etki & O$-$C şekli & RV'de & Örtülmelerde & Basit test\\
\midrule
1 Gelgit & Aşağı parabol, sabit $\dot P$ & Eğilim yok & Geçişle aynı yönde & Q$'_\star\sim10^5$–$10^7$ mi?\\
2 Apsis & Sinüs (kısa veride parabol) & Eğilim yok & Geçişle \textbf{ters} yönde & Örtülme zamanları; $e$\\
3 Rømer & Uzun sinüs; kısa veride parabol (±) & \textbf{Eğilim var}, $\dot P = P\dot v_r/c$ & Aynı yönde & Arşiv RV eğimi\\
4 Applegate & Yarı periyodik (5–50 yıl) & Yok & Aynı yönde & Etkin G/K yıldız? Periyodogram\\
5 3. cisim (dinamik) & Kısa dönemli TTV / zorlanmış presesyon & Ek gezegen sinyali olabilir & Duruma göre & Periyodogram; literatür\\
6 Kütle kaybı & Çok yavaş, tekdüze (çoğunlukla +) & Yok & Aynı yönde & $a/a_\text{Roche}\approx1$ mi?\\
Veri hatası & Tek nokta / tek kaynak & -- & -- & Zaman sistemi kontrolü; şüpheli kaynakla/kaynaksız uydurma\\
\bottomrule
\end{tabularx}
\end{center}

% ===================================================== BÖLÜM 8
\section{Gelgit mi, Rømer mi? Radyal Hız (RV) Adımı}
\label{sec:rv}

\subsection{Duyduğun şey doğru mu? EVET.}
“Sadece O$-$C'ye bakarak gelgit bozulmasını üçüncü cismin ivmesinden (Rømer) matematiksel olarak ayıramayız” sözü
\textbf{doğrudur} (Bouma vd. 2019, 2020; Winn \& Stefánsson 2025). İspatı üç satır:

Sistem bize göre sabit bir $\dot v_r$ ivmesiyle hareket etsin. $t$ anında sistemin uzaklığı
$d(t) = d_0 + v_0 t + \tfrac12 \dot v_r t^2$. Işık bize $d/c$ sürede gelir. $E$'inci geçiş $t_E \approx T_0 + PE$
anında olur ve bize
\[
T_E = T_0 + PE + \frac{v_0 P E}{c} + \frac{1}{2}\,\frac{\dot v_r P^2}{c}\,E^2
\]
anında ulaşır. Bu ifade, \textbf{kuadratik efemerisle birebir aynı biçimdedir}: $T_E = T_0 + P'E + \tfrac12 (dP/dE)E^2$,
burada $dP/dE = \dot v_r P^2/c$, yani $\dot P = P\dot v_r/c$. O$-$C verisi yalnızca “$E^2$ teriminin katsayısını”
ölçer; bu katsayının gelgitten mi ivmeden mi geldiğini \emph{aynı veriden} anlamak imkânsızdır.

\textbf{İstisnalar} (zamanlamanın kendisinin bir şey söyleyebildiği durumlar):
(1) Eşlikçinin dönemi gözlem süresiyle karşılaştırılabilirse ivme değişir; kübik terim ya da sinüs gerekir
(WASP-4'te kübik test $\Delta$BIC = 13 verdi, Bölüm~\ref{sec:karar}).
(2) Örtülme zamanları varsa apsis hareketi ayrılır (ama gelgit ile Rømer yine ayrılmaz: ikisinde de örtülmeler aynı yönde kayar).

\subsection{Çözüm: arşivden RV}
Gelgit sistemin kütle merkezini \emph{ivmelendirmez}; Rømer etkisi ise tanımı gereği ivmedir. RV ölçümleri ivmeyi
\textbf{doğrudan} verir. Bu yüzden:
\[
\text{Test: } \quad \dot P_\text{gözlenen} \overset{?}{=} \dot P_\text{Rømer} = \frac{P\,\dot v_r}{c} .
\]
\begin{itemize}[leftmargin=*]
\item $|\dot P_\text{gözlenen} - \dot P_\text{Rømer}| < 2\sigma$: \textbf{Rømer tek başına açıklayabilir} (gelgite gerek yok).
\item Fark $> 3\sigma$ ve $\dot P_\text{Rømer} \approx 0$: \textbf{Rømer dışlanır}; gelgit (ya da diğerleri) gerekir.
\item Arada: “kısmen”. Kalan $\dot P_\text{gözlenen} - \dot P_\text{Rømer}$ diğer etkilere bırakılır.
\end{itemize}

\subsection{Hangi RV verisi, nereden? (tayfçeker gerekmez)}
\textbf{Önemli ayrım:} Doğrulanmış her sıcak Jüpiterin RV verisi \emph{vardır} (gezegenin kütlesi RV'den ölçülür). Ama
çoğu yalnızca keşif döneminde, birkaç ay içinde alınmıştır. Rømer testi için gereken, \textbf{yıllara yayılmış} RV'dir:
yıldızın hızının yıllar içinde düzenli değişip değişmediği ancak öyle görülür. Ayrıca farklı spektrograflar farklı sıfır
noktası verir; her aletin sıfırı ayrı uydurulur ve bu, eğim bilgisinin bir kısmını yer. Bu yüzden asıl soru “RV var mı?”
değil, “bu RV'lerle Rømer öngörüsü ne kadar hassas çıkıyor?” sorusudur (Bölüm~\ref{sec:secim}, K3).

\begin{center}
\small
\begin{tabularx}{\textwidth}{l >{\raggedright\arraybackslash}X >{\raggedright\arraybackslash}X}
\toprule
Kaynak & İçerik & Erişim\\
\midrule
\textbf{DACE} (Cenevre) & CORALIE, HARPS, HARPS-N, SOPHIE, HIRES, ESPRESSO arşiv RV'leri. 14 adayın 12'si var
  (HAT-P-37 ve KELT-16 yok) & \texttt{dace.unige.ch}; Python: \texttt{dace\_rv\_indir.py} (hesap gerekmez)\\
Knutson vd. (2014) & 51 sıcak Jüpiterin Keck/HIRES RV'leri; HAT-P-7 (44), WASP-12 (30), WASP-4 (6), WASP-10, WASP-22, WASP-1, XO-4 &
  VizieR \texttt{J/ApJ/785/126}, \texttt{table3}\\
Bouma vd. (2020) & WASP-4 için 48 RV \textbf{ve Tablo 5: pek çok sıcak Jüpiter için RV eğiminden öngörülen $\dot P$ hazır} &
  VizieR \texttt{J/ApJ/893/L29}, \texttt{table2}, \texttt{table5}\\
Bonomo vd. (2017) & 231 geçiş yapan gezegen; HARPS-N RV'leri (TrES-1 için 15 ölçüm) ve her gezegen için RV eğimi &
  VizieR \texttt{J/A+A/602/A107}, \texttt{table1} (RV), \texttt{table5} (özet), \texttt{table8} (eğim)\\
Yee vd. (2020) & WASP-12 için 72 RV ve 19 örtülme zamanı & VizieR \texttt{J/ApJ/888/L5}, \texttt{table4}, \texttt{table5}\\
HARPS RVBank (Trifonov vd. 2020) & Yeniden işlenmiş HARPS RV'leri; adaylardan yalnız WASP-4 ve WASP-22 var & VizieR \texttt{J/A+A/636/A74}\\
\bottomrule
\end{tabularx}
\end{center}
\textbf{Bouma vd. (2020) Tablo 5}'te birim m/s/gün'dür (VizieR başlığında m/s/yr yazsa da değerler m/s/gün ile tutarlıdır;
ör. WASP-4: $-0{,}0422$). Hazır değerler: WASP-12\,b $\dot P_\text{Rømer} = -0{,}10\pm1{,}07$; WASP-4\,b $-5{,}94\pm0{,}39$;
HAT-P-7\,b $+14{,}99\pm0{,}88$\msyil.

\textbf{DACE'den indirmek} (kod Bölüm~\ref{sec:kod}):
\begin{lstlisting}[style=kabuk]
pip install dace-query
python dace_rv_indir.py WASP-12 1.091418901 2457607.519305       # rv_WASP-12.csv olusur
WASP-12: 90 RV, 13.9 yil, aletler: {'ESPRESSO19': 28, 'SOPHIE': 62}
python rv_analiz.py rv_WASP-12.csv WASP-12 1.091418901 2457607.519305 -31.05 0.82
\end{lstlisting}

\subsection{Sonuçlar (kendi hesabım, \texttt{rv\_analiz.py})}
\begin{center}
\small
\begin{tabularx}{\textwidth}{l >{\raggedright\arraybackslash}X r r r >{\raggedright\arraybackslash}X}
\toprule
Hedef & RV verisi & $\dot v_r$ (m/s/yıl) & $\dot P_\text{Rømer}$ & $\dot P_\text{O$-$C}$ & Sonuç\\
\midrule
WASP-12\,b & DACE, 90 RV, 13,9 yıl & $-2{,}5\pm5{,}4$ & $-0{,}8\pm1{,}7$ & $-31{,}05\pm0{,}82$ & Fark 16$\sigma$: \textbf{Rømer dışlandı}\\
WASP-4\,b  & Bouma, 48 RV, 11,9 yıl & $-15{,}44\pm0{,}83$ & $-5{,}95\pm0{,}32$ & $-5{,}11\pm0{,}30$ & Fark 1,9$\sigma$: \textbf{Rømer açıklıyor}\\
HAT-P-7\,b & Knutson, 43 RV, 5,9 yıl & $+20{,}6\pm1{,}9$ & $+13{,}1\pm1{,}2$ & $+28{,}14\pm3{,}26$ & Fark 4,3$\sigma$: \textbf{kısmen}\\
TrES-1\,b  & Bonomo, 15 RV, 2,8 yıl & $-2{,}9\pm2{,}5$ & $-2{,}5\pm2{,}2$ & $-12{,}70\pm1{,}90$ & Fark 3,5$\sigma$: \textbf{Rømer dışlandı}\\
\bottomrule
\end{tabularx}
\end{center}
Sağlama: WASP-4 için bulduğum $-15{,}44$ m/s/yıl $= -0{,}0423$ m/s/gün, Bouma vd.'nin $-0{,}0422\pm0{,}0028$ değeriyle
aynı. HAT-P-7 için $+20{,}6$, Winn vd. (2009)'un $21{,}5\pm2{,}6$ değeriyle uyumlu. WASP-12 için Knutson'ın 29 RV'siyle de
aynı sonuç çıkıyor ($-3{,}6\pm1{,}6$; Rømer yine dışlanıyor).

\begin{dikkat}[Hangi RV setini kullandığın sonucu değiştirebilir]
WASP-4 için yalnız DACE verisiyle (CORALIE + HARPS) Rømer öngörüsü $+0{,}5\pm1{,}4$ çıkıyor; Bouma'nın setiyle
$-5{,}95\pm0{,}32$. Fark, Bouma vd.'nin 2019'da aldığı ve DACE'de olmayan HIRES ölçümlerinden geliyor: zaman tabanını
uzatan onlar. Ders: \textbf{bir hedefin bütün yayımlanmış RV'lerini birleştir}, sonucun veri setine duyarlılığını raporla
(bu, İP4'ün bir parçası).
\end{dikkat}

\textbf{HAT-P-7'de neden “kısmen”?} İki bilinen sebep: (1) Knutson vd. (2014), HAT-P-7'nin RV eğiliminin tepe yapıp
döndüğünü bildirir; yani ivme sabit değil ve 2008–2013'teki eğim, 2009–2025 boyunca ortalama ivmeyi temsil etmeyebilir
(LTTE'nin eğriliği). (2) O$-$C'deki 9 eski zaman “BJD” olarak etiketli; zaman sistemi doğrulanmadan $\dot P_\text{O$-$C}$
kesinleşmez (Bölüm~\ref{sec:kavram}). Proje sırasında ikisi de kontrol edilecek. Bu, projenin gerçek bir araştırma sorusudur.

\textbf{TrES-1 ne anlatıyor?} RV'ler yıldızın 2,8 yılda belirgin bir ivme göstermediğini söylüyor: Rømer olsaydı RV
yılda yaklaşık $-14{,}5$ m/s değişmeliydi (aşağıdaki hesap), HARPS-N'in $\sim$5 m/s hassasiyetiyle bu kolayca görülürdü.
Gelgit de Q$'_\star \approx 8\times10^2$ gerektiriyor. Geriye apsis hareketi (Hagey vd. 2025'in önerisi), dinamik 3. cisim ve
Applegate kalıyor; projenin en ilginç sorusu burada.

\subsection{B Planı: Yeterli RV verisi olmazsa ne yaparım?}
Seçtiğimiz dört hedefte yeterli RV var; ama proje sırasında bir hedefin RV'si yetersiz kalırsa (ya da yedek hedefe geçilirse)
izlenecek yol:
\begin{enumerate}[leftmargin=*]
\item \textbf{Literatürü tara}: Hedefe özel makalelerde RV eğimi yayımlanmış mı? Yayımlanmış eğimi al,
      $\dot P_\text{Rømer} = P\dot v_r/c$ ile hesapla. Tayfçeker gerekmez; makaledeki sayıyı kullanırsın.
\item \textbf{Ters soru sor (test edilebilir öngörü)}: “Bu dönem değişimi tamamen Rømer olsaydı yıldızın RV'si ne hızla
      değişmeliydi?” $\dot v_r = c\,\dot P/P$. Örnek (TrES-1): $\dot v_r = 2{,}998\times10^8 \times
      (-12{,}7\times10^{-3}/3{,}156\times10^7) / (3{,}030\times86400) \approx -0{,}040$ m/s/gün $\approx -14{,}5$ m/s/yıl.
      Bu sayı raporlanır: gelecekte RV ölçen herkes bunu doğrudan sınayabilir.
\item \textbf{RV gerektirmeyen testleri uygula} (\texttt{fizik\_testleri.py}): gereken Q$'_\star$, Roche oranı, kübik/periyodogram,
      yıldız etkinliği.
\item \textbf{Hükmü dürüstçe ver}: Ör. “Gelgit fiziksel olarak olası değil; kütle kaybı dışlandı; Rømer ve Applegate RV
      olmadan ayrılamadı; Rømer için gereken RV eğimi şudur.” Bu da 2209-A için geçerli bir sonuçtur.
\end{enumerate}

% ===================================================== BÖLÜM 9
\section{Karar Ağacı: Hangi Etki?}
\label{sec:karar}

Her hedef için aşağıdaki adımları \textbf{sırayla} uygula. Eşikler önceden yazılmıştır; sonuca bakıp değiştirme
(değiştirirsen eski ve yeni sürümün sonuçlarını birlikte raporla).

\begin{enumerate}[label=\textbf{A\arabic*.},leftmargin=*]
\setcounter{enumi}{-1}
\item \textbf{Veri temizliği}: Bütün zamanlar BJD$_\mathrm{TDB}$ mi? Epoklar tam sayıya yakın mı ($|x-E|<0{,}02$)?
      Aynı geçiş iki kez var mı? Belirsiz kaynaklar ile ve onlarsız iki uydurma yapıldı mı?
\item \textbf{Değişim var mı?} $|\dot P|/\sigma \ge 3$ \emph{ve} $\Delta$BIC $>10$ \emph{ve}
      tek bir kaynağı çıkarmak sonucu yok etmiyor. Hayır $\Rightarrow$ \textbf{“anlamlı değişim yok”} (dur).
\item \textbf{Rømer testi} (RV varsa): $|\dot P - \dot P_\text{Rømer}| < 2\sigma$ $\Rightarrow$ \textbf{Rømer/3. cisim
      (LTTE)}. Fark $>3\sigma$ ve $\dot P_\text{Rømer}\approx0$ $\Rightarrow$ Rømer dışlandı, A3'e geç.
      Arada $\Rightarrow$ “kısmen”; kalan $\dot P - \dot P_\text{Rømer}$ ile A3'e geç. RV yetersizse B planı (Bölüm~\ref{sec:rv}).
\item \textbf{Dönem uzuyor mu?} Evet $\Rightarrow$ gelgit dışlanır (yıldız yavaş dönüyor); A5–A7'ye bak.
\item \textbf{Gelgit olabilirliği}: Gereken Q$'_\star \ge 10^4$ ve kübik terim gereksiz ($\Delta$BIC$_\text{kübik} < 10$)
      ve (varsa) örtülmeler aynı yönde $\Rightarrow$ \textbf{gelgit en olası}. Q$'_\star < 10^4$ $\Rightarrow$ gelgit
      “fiziksel olarak olası değil”.
\item \textbf{Apsis}: $e$ üst sınırından genlik gözlenen sinyali açıklayabilir mi? Örtülmeler ters yönde mi?
      Kübik/sinüs tercih ediliyor mu?
\item \textbf{Applegate / dinamik 3. cisim}: Etkin G/K yıldızı mı? Periyodogramda (alias olmayan, FAP $<0{,}01$) tepe
      var mı? Literatürde ek gezegen/eşlikçi var mı?
\item \textbf{Kütle kaybı}: $a/a_\text{Roche} < 1{,}5$ mı? Değilse dışla.
\end{enumerate}
Her etki için hüküm: \textbf{dışlandı / olası değil / olası / en olası / ayrılamadı}.

\subsection{Ön sonuçlarla doldurulmuş karar tablosu}
\begin{center}
\small
\begin{tabularx}{\textwidth}{l >{\raggedright\arraybackslash}X >{\raggedright\arraybackslash}X >{\raggedright\arraybackslash}X >{\raggedright\arraybackslash}X}
\toprule
Etki & WASP-12\,b & WASP-4\,b & HAT-P-7\,b & TrES-1\,b\\
\midrule
Değişim var mı? & Evet ($-31{,}1$) & Evet ($-5{,}1$) & Evet ($+28{,}1$) & Evet ($-12{,}7$)\\
1 Gelgit & \textbf{En olası} (Q$'_\star=1{,}9\times10^5$; kübik gereksiz, $\Delta$BIC 8,9) & Olası ama zorlayıcı (Q$'_\star=7{,}7\times10^4$) & Dışlandı (dönem uzuyor) & Olası değil (Q$'_\star=8\times10^2$)\\
2 Apsis & Dışlandı (örtülmeler aynı yönde) & Ayrılamadı (örtülme yok) & Ayrılamadı & \textbf{En olası aday} (Hagey vd. 2025)\\
3 Rømer & Dışlandı (16$\sigma$, DACE) & \textbf{En olası} (1,9$\sigma$; $\dot P$ zamanla değişiyor: kübik $\Delta$BIC 13) & \textbf{Kısmen} (eğim 13'ü açıklıyor) & Dışlandı (3,5$\sigma$, HARPS-N)\\
4 Applegate & Olası değil (F yıldızı) & Dışlanamadı (G7V) & Olası değil (F yıldızı) & Dışlanamadı (K0V)\\
5 3. cisim (din.) & Kanıt yok & Kanıt yok (eşlikçi uzakta) & Bilinen eşlikçiler (RV + görüntüleme) & Önerilmiş yakın eşlikçi (Hagey vd.)\\
6 Kütle kaybı & Roche'a yakın (1,18) ama kısalmayı açıklayamaz & İhmal (1,76) & İhmal (2,54) & İhmal (3,09)\\
\bottomrule
\end{tabularx}
\end{center}
Bu tablo projenin \emph{başlangıç} durumudur. Proje boyunca her hücre ayrıntılı testlerle kesinleştirilecek; asıl
sonuç bu tablonun son halidir.

